\section{Introduction}
\label{sec:introduction}

Simulation studies are routinely used to compare scheduling policies when
exhaustive real-system experimentation is too expensive, slow, or difficult
to control. Such studies do not produce scheduler conclusions in the
abstract: a comparative ordering is produced under a chosen workload,
operating region, resource model, and evaluation metric. The same simulation
campaign can therefore answer ``which scheduler is better'' only relative to
the conditions under which that comparison was made. If the workload source,
load region, metric, or simulator fidelity boundary changes, the comparative
ranking may remain stable, degrade gradually, or reverse. This paper studies
that portability problem as the primary output of a trace-driven simulation
experiment, rather than treating it as a secondary sensitivity check after a
single scheduler leaderboard has already been reported.

LLM serving is a useful and practically important setting for this question.
Modern LLM serving targets high-performance, GPU-resident inference systems
whose scheduling decisions govern how requests share costly accelerator,
memory, and batching resources, and serving stacks have grown increasingly
sophisticated in how they make those decisions. Beyond simple
first-come-first-served execution, modern serving stacks employ continuous
batching and paged key-value (KV) cache management
\citep{kwon2023vllm}, iteration-level scheduling \citep{yu2022orca}, chunked
prefill to balance prefill and decode work within a batch
\citep{agrawal2024sarathi}, disaggregated prefill/decode execution across
dedicated hardware pools \citep{zhong2024distserve,patel2024splitwise},
KV-cache-centric disaggregated architectures \citep{qin2025mooncake},
dynamic instance-level rebalancing \citep{sun2024llumnix}, explicit
fairness objectives \citep{sheng2024vtc}, output-length-aware admission and
memory packing \citep{jin2023s3}, and learning-to-rank decode-priority
policies \citep{efficientllmscheduling2024}. Each of these mechanisms was
introduced, and evaluated, against a particular choice of simulation or
experimental conditions: a specific trace, a specific arrival process, a
specific mix of prompt and output lengths, a specific operating load, and a
specific metric.

A separate line of work has documented that these workloads are themselves
far from uniform. Public production traces from different services and time
periods differ substantially in burstiness, request concurrency, prompt and
output length distributions, and failure characteristics
\citep{burstgpt2025,patel2024splitwise,stojkovic2025dynamollm}; large-scale
characterization studies of KV-cache reuse and of year-long production
traffic report further axes of heterogeneity across clients, models, and
time \citep{kvcache_wild2025,ayearinllmserving2026}; and dedicated
workload-generation frameworks exist precisely because no single fixed trace
is thought to be representative of production traffic in general
\citep{servegen2026}. Individually, prior work has therefore already
established two things: that LLM-serving mechanisms are numerous and
mechanistically diverse, and that the workloads used to evaluate them are
themselves heterogeneous and shift across source, time, and operating
region.

What has not been directly asked for LLM-serving scheduler comparisons is
the simulation-comparison question that follows from combining these two
observations: whether comparative scheduler conclusions transfer across
contexts. The typical evaluation question in this literature is: \emph{how
well does scheduler $A$ perform, relative to scheduler $B$, on the
workload(s) chosen for this paper's experiments?} Because a simulation study
necessarily uses only a finite, chosen set of heterogeneous workloads, a
comparative claim of this form is only as useful as its
\emph{portability}: does the same comparative ordering of schedulers hold if
the workload distribution, the operating load, or the evaluation metric is
changed to another one a practitioner might reasonably care about? This
question differs both from workload heterogeneity itself and from any single
scheduler's sensitivity to load or service-level objectives (SLOs): it
concerns the reliability of the \emph{comparison}, not of any individual
system. LSSP is designed to measure this transferability directly under a
fixed simulation protocol.

We formalize this as the portability of a comparative scheduler ranking
$R(W, M, L)$, where $W$ denotes a workload distribution, window, or source;
$M$ denotes an evaluation metric; and $L$ denotes an operating or load
region. $R(W, M, L)$ is the ordering (full or partial) that a fixed panel of
scheduling policies induces as the comparative output of a simulation
experiment under a fixed protocol, holding $W$, $M$, and $L$ at specified
values. Ranking portability asks how $R$ behaves as $W$, $M$, or $L$ is
varied: whether the ordering of policies is preserved, where and how it
reverses, and how much evidence (how many independent workload windows,
drawn from how many sources) is required before a reported ranking can be
trusted to generalize beyond the exact windows used to produce it.

This question matters because a scheduler comparison is scientifically and
practically useful only insofar as its comparative conclusion transfers to
the workload distributions a reader actually cares about. A benchmark result
of the form ``scheduler $A$ outperforms scheduler $B$'' is not a claim about
$A$ and $B$ in the abstract; it is implicitly a claim about $A$ and $B$
\emph{under the workload, metric, and load conditions the benchmark used}.
If that comparative ordering is fragile under source, metric, or
load-region substitution, then benchmark results risk being read as more
general than they are, and evaluation effort risks being spent on a single
workload family when the resulting conclusion may not transfer. Conversely,
if comparative orderings are portable across independent sources, metrics,
and load regions, that is itself a useful, non-obvious finding: it would
mean a leaner benchmark protocol is sufficient for future scheduler
comparisons, and it would give scheduler designers a much stronger warrant
for claims of general improvement.

We refer to the benchmark built around this question as the LLM-Serving
Scheduler Portability Benchmark (LSSP). LSSP asks the following research
questions, fixed under a preregistered protocol before any comparative
outcome was generated (\S\ref{sec:study-design}):

\begin{description}
  \item[RQ1.] How stable are scheduler rankings across independent
    workload sources?
  \item[RQ2.] How stable are scheduler rankings under temporal, provider,
    and domain shifts?
  \item[RQ3.] To what extent do rankings obtained on synthetic stress
    workloads transfer to rankings on independent real-trace-derived
    workloads?
  \item[RQ4.] Which source-native observable workload characteristics are
    associated with cross-distribution scheduler rank reversals?
  \item[RQ5.] How many independent workload windows are required before a
    comparative scheduler ranking becomes statistically stable?
  \item[RQ6.] Do representative simulated scheduler rankings and
    cross-workload rank reversals reproduce on a real serving engine?
\end{description}

Load-level dependence is deliberately treated as a secondary
robustness/sensitivity axis applied to RQ1--RQ4, rather than as a headline
research question in its own right, to avoid duplicating a closely related
prior load-scaling study on an overlapping window set
(\S\ref{sec:study-design}).

\paragraph{Contributions} This study executes that preregistered protocol
and reports its structurally validated results. This work contributes:

\begin{itemize}
  \item a formulation of comparative scheduler-ranking portability,
    $R(W, M, L)$, as an object of simulation-based selection and comparison
    studies, distinct from both workload characterization and
    single-scheduler evaluation;
  \item a preregistered paired experimental design for comparing a fixed
    panel of scheduling alternatives across independently sourced workload
    windows, calibrated operating regions, and metrics, so that
    ranking-portability statistics are computed on matched observations
    rather than separately tuned per-source experiments;
  \item a deterministic, trace-driven discrete-event simulation execution
    model for LLM-serving scheduler comparisons, with repeated executions
    used as verification of bit-identical outcomes rather than as
    stochastic replicates (\S\ref{sec:execution-model});
  \item an uncertainty-aware reversal-detection methodology that
    distinguishes practically meaningful rank reversals from microscopic
    sign changes, using preregistered practical-margin and
    confidence-interval criteria, and that treats every
    completion-conditioned metric's undefined cases explicitly rather than
    by imputation (\S\ref{sec:methodology});
  \item a benchmark sample-complexity analysis quantifying how many
    independent workload windows are required before a comparative ranking
    reproduces, applied to this study's own frozen 120-window corpus
    (\S\ref{sec:sample-complexity});
  \item a characterization of the workload corpus underlying the
    simulation study, establishing that the three sources are measurably
    different in observed descriptor distributions and strongly separable
    under a leakage-resistant, chronologically blocked evaluation
    (\S\ref{sec:workloads});
  \item a selected-case physical validation on real vLLM/GPU hardware
    (RQ6, \S\ref{sec:real-system}) used explicitly to characterize
    simulator fidelity and domain of validity for the benchmark's largest
    identified simulated reversal and a paired stable-control comparison;
    and
  \item a public artifact release of the benchmark's frozen workload
    identities, policy panel, protocol, simulation outputs, and analysis
    artifacts, to support independent reproduction (\S\ref{sec:artifact}).
\end{itemize}

Comparative rankings prove neither uniformly stable nor uniformly
unstable, but source-pair- and load-region-dependent
(Sections~\ref{sec:results}--\ref{sec:real-system}); we do not claim that
any scheduler in the panel wins in general (\S\ref{sec:discussion}).
